\subsection{Los $\mathscr{P}_{X/S}^n(\mathscr{F})$.}
\label{IV.16.7}

\begin{env}[16.7.1]
\label{IV.16.7.1}
Sean $f:X\to S$ un morfismo de preesquemas, $\mathscr{F}$ un
$\mathcal{O}_X$-módulo. Denotemos por $X_{\Delta_f}^{(n)}$ la \emph{$n$-ésima
vecindad infinitesimal de $X$} para el morfismo diagonal
\oldpage[IV-4]{37}
$\Delta_f:X\to X\times_S X$, por
$h_n:X_{\Delta_f}^{(n)}\to X\times_S X$ el morfismo canónico (16.1.2), y
consideremos los dos morfismos compuestos
\[
  p_1^{(n)}:X_{\Delta_f}^{(n)}\xrightarrow{h_n}X\times_S X
  \xrightarrow{p_1}X,\qquad
  p_2^{(n)}:X_{\Delta_f}^{(n)}\xrightarrow{h_n}X\times_S X
  \xrightarrow{p_2}X
\]
de modo que, por definición, $p_1^{(n)}$ corresponde al homomorfismo de
haces de anillos $\mathcal{O}_X\to\mathscr{P}_{X/S}^n$ que hemos
elegido para definir la estructura de $\mathcal{O}_X$-álgebra sobre
$\mathscr{P}_{X/S}^n$ (16.3.5), y $p_2^{(n)}$ al homomorfismo de haces
de anillos $d_{X/S}^n:\mathcal{O}_X\to\mathscr{P}_{X/S}^n$ (16.3.6). Como
$X_{\Delta_f}^{(n)}$ y $X$ tienen el mismo espacio subyacente, se puede escribir
\[
  \mathscr{P}_{X/S}^n
  =(p_1^{(n)})_*((p_2^{(n)})^*(\mathcal{O}_X)).
  \tag{16.7.1.1}
  \label{IV.16.7.1.1}
\]
Más generalmente, pondremos
\[
  \mathscr{P}_{X/S}^n(\mathscr{F})
  =(p_1^{(n)})_*((p_2^{(n)})^*(\mathscr{F}))
  \tag{16.7.1.2}
  \label{IV.16.7.1.2}
\]
de modo que
$\mathscr{P}_{X/S}^n=\mathscr{P}_{X/S}^n(\mathcal{O}_X)$; por definición,
$\mathscr{P}_{X/S}^n(\mathscr{F})$ es, por tanto, un $\mathcal{O}_X$-módulo.
\end{env}

\begin{env}[16.7.2]
\label{IV.16.7.2}
Si se vuelve a las definiciones de las imágenes inversas de módulos sobre los
espacios anillados (0, 4.3.1) y se tiene en cuenta que
$X_{\Delta_f}^{(n)}$ y $X$ tienen el mismo espacio subyacente, se ve que se puede
escribir también la definición (16.7.1.2) en la forma
\[
  \mathscr{P}_{X/S}^n(\mathscr{F})
  =\mathscr{P}_{X/S}^n\otimes_{\mathcal{O}_X}\mathscr{F}
  \tag{16.7.2.1}
  \label{IV.16.7.2.1}
\]
pero debe tenerse cuidado de que, al interpretar el signo $\otimes$,
$\mathscr{P}_{X/S}^n$ está dotado de su estructura de $\mathcal{O}_X$-módulo
definida por el homomorfismo de haces de anillos
$d_{X/S}^n:\mathcal{O}_X\to\mathscr{P}_{X/S}^n$. Resulta inmediatamente de esta
fórmula (o directamente de (16.7.1.2)) que
$\mathscr{P}_{X/S}^n(\mathscr{F})$ está canónicamente dotado de una estructura de
$\mathscr{P}_{X/S}^n$-módulo.
\end{env}

\begin{proposition}[16.7.3]
\label{IV.16.7.3}
\begin{enumerate}
\item[(i)] El funtor
$\mathscr{F}\mapsto\mathscr{P}_{X/S}^n(\mathscr{F})$ de la categoría de los
$\mathcal{O}_X$-módulos en la de los $\mathscr{P}_{X/S}^n$-módulos es
exacto por la derecha y conmuta con los límites inductivos cualesquiera; es exacto
cuando $\mathscr{P}_{X/S}^n$ es un $\mathcal{O}_X$-módulo plano.
\item[(ii)] Si $\mathscr{F}$ es un $\mathcal{O}_X$-módulo cuasicoherente
(resp. de tipo finito, resp. de presentación finita),
$\mathscr{P}_{X/S}^n(\mathscr{F})$ es un $\mathscr{P}_{X/S}^n$-módulo
cuasicoherente (resp. de tipo finito, resp. de presentación finita).
\end{enumerate}
\end{proposition}

Las afirmaciones de (i) resultan inmediatamente de la fórmula (16.7.2.1) y de la
consideración de la simetría de $\mathscr{P}_{X/S}^n$ (16.3.4). Las afirmaciones
de (ii) se deducen de la exactitud por la derecha del funtor
$\mathscr{F}\mapsto\mathscr{P}_{X/S}^n(\mathscr{F})$.

\begin{env}[16.7.4]
\label{IV.16.7.4}
Las dos estructuras de $\mathcal{O}_X$-módulo de $\mathscr{P}_{X/S}^n$
definen sobre $\mathscr{P}_{X/S}^n(\mathscr{F})$ dos estructuras de
$\mathcal{O}_X$-módulo que, por lo demás, conmutan y, por tanto, una estructura de
$\mathcal{O}_X$-bimódulo. Conviene llamar \emph{izquierda} a la estructura
proveniente del homomorfismo estructural
$\mathcal{O}_X\to\mathscr{P}_{X/S}^n$ (elegido en (16.3.5)) y
\emph{derecha} a la proveniente del homomorfismo
$d_{X/S}^n:\mathcal{O}_X\to\mathscr{P}_{X/S}^n$. Es decir, para todo
abierto $U$ de $X$ y toda terna de elementos
$a\in\Gamma(U,\mathcal{O}_X)$, $b\in\Gamma(U,\mathscr{P}_{X/S}^n)$,
$t\in\Gamma(U,\mathscr{F})$, se tiene por definición
\[
  a(b\otimes t)=(ab)\otimes t,\qquad
  (b\otimes t)a=(b.d^na)\otimes t=b\otimes(at)
  =(d^na).(b\otimes t).
  \tag{16.7.4.1}
  \label{IV.16.7.4.1}
\]
La estructura de $\mathcal{O}_X$-módulo proveniente de la definición (16.7.1.2)
es, por tanto, con estas convenciones, la estructura de $\mathcal{O}_X$-módulo a la
izquierda.

\oldpage[IV-4]{38}
Si $\mathscr{F}$ es un $\mathcal{O}_X$-módulo cuasicoherente, lo mismo ocurre
con $\mathscr{P}_{X/S}^n(\mathscr{F})$ para cualquiera de sus
estructuras de $\mathcal{O}_X$-módulo. Si además $\mathscr{F}$ es de tipo
finito (resp. de presentación finita) y $f:X\to S$ es localmente de tipo finito
(resp. localmente de presentación finita),
$\mathscr{P}_{X/S}^n(\mathscr{F})$ es (para cualquiera de sus
estructuras de $\mathcal{O}_X$-módulo) de tipo finito (resp. de presentación
finita), como resulta de (16.3.9) y de (16.4.22).
\end{env}

\begin{env}[16.7.5]
\label{IV.16.7.5}
La definición (16.7.2.1) implica la existencia de un homomorfismo de haces
de grupos conmutativos
\[
  d_{X/S,\mathscr{F}}^n:\mathscr{F}
  \longrightarrow\mathscr{P}_{X/S}^n(\mathscr{F})
  \qquad\text{(también denotado por $d_{X/S}^n$)}
  \tag{16.7.5.1}
  \label{IV.16.7.5.1}
\]
tal que, con las notaciones de (16.7.4), se tenga
\[
  d_{X/S,\mathscr{F}}^n(t)=1\otimes t
  \tag{16.7.5.2}
  \label{IV.16.7.5.2}
\]
y, por consiguiente, en virtud de (16.7.4.1),
\[
  d_{X/S,\mathscr{F}}^n(at)
  =(1\otimes t)a=(d_{X/S,\mathscr{F}}^n(t)).a
  \tag{16.7.5.3}
  \label{IV.16.7.5.3}
\]
\[
  d_{X/S,\mathscr{F}}^n(at)
  =(d_{X/S}^n(a)).(1\otimes t)
  =(d_{X/S}^n(a))(d_{X/S,\mathscr{F}}^n(t)).
  \tag{16.7.5.4}
  \label{IV.16.7.5.4}
\]
Es, por tanto, $\mathcal{O}_X$-lineal para la estructura de
$\mathcal{O}_X$-módulo a la derecha sobre $\mathscr{P}_{X/S}^n(\mathscr{F})$, y
\emph{semilineal} (respecto del automorfismo $\sigma$ (16.3.4)) para
la estructura de $\mathcal{O}_X$-módulo a la izquierda.
\end{env}

\begin{proposition}[16.7.6]
\label{IV.16.7.6}
El $\mathcal{O}_X$-módulo a la izquierda $\mathscr{P}_{X/S}^n(\mathscr{F})$ está
engendrado por la imagen de $\mathscr{F}$ por el homomorfismo canónico
$d_{X/S,\mathscr{F}}^n$.
\end{proposition}

Esto resulta inmediatamente de (16.7.5.3) y del caso particular correspondiente a
$\mathscr{F}=\mathcal{O}_X$ (16.3.8).

\begin{env}[16.7.7]
\label{IV.16.7.7}
Los homomorfismos canónicos de haces de anillos
\[
  \varphi_{nm}:\mathscr{P}_{X/S}^m\longrightarrow\mathscr{P}_{X/S}^n
\]
para $n\leq m$ (16.1.2) definen, en virtud de (16.7.2.1),
homomorfismos canónicos
\[
  \mathscr{P}_{X/S}^m(\mathscr{F})
  \longrightarrow\mathscr{P}_{X/S}^n(\mathscr{F})\qquad(n\leq m)
\]
que son homomorfismos de $\mathcal{O}_X$-bimódulos en virtud de (16.1.6)
y (16.7.4.1); además, se tienen diagramas conmutativos
\[
  \xymatrix{
    \mathscr{P}_{X/S}^m(\mathscr{F})\ar[rr]
      &&\mathscr{P}_{X/S}^n(\mathscr{F})\\
    &\mathscr{F}\ar[ul]^-{d_{X/S,\mathscr{F}}^m}
      \ar[ur]_-{d_{X/S,\mathscr{F}}^n}
  }
\]
Se tiene así un sistema proyectivo de $\mathcal{O}_X$-bimódulos
$(\mathscr{P}_{X/S}^n(\mathscr{F}))$, y se pone
\[
  \mathscr{P}_{X/S}^{\infty}(\mathscr{F})
  =\varprojlim\mathscr{P}_{X/S}^n(\mathscr{F}).
  \tag{16.7.7.1}
  \label{IV.16.7.7.1}
\]
Además, lo que precede muestra que los homomorfismos (16.7.5.1) forman un
sistema proyectivo de homomorfismos y definen, por tanto, un homomorfismo
canónico
\[
  d_{X/S,\mathscr{F}}^{\infty}:\mathscr{F}
  \longrightarrow\mathscr{P}_{X/S}^{\infty}(\mathscr{F}).
  \tag{16.7.7.2}
  \label{IV.16.7.7.2}
\]
\end{env}

\oldpage[IV-4]{39}

\begin{env}[16.7.8]
\label{IV.16.7.8}
Sean $\mathscr{F}$, $\mathscr{G}$ dos $\mathcal{O}_X$-módulos; resulta
inmediatamente de la definición (16.7.2.1) que se tiene un isomorfismo canónico de
$\mathscr{P}_{X/S}^n$-módulos
\[
  \mathscr{P}_{X/S}^n(\mathscr{F}\otimes_{\mathcal{O}_X}\mathscr{G})
  \xrightarrow{\sim}
  \mathscr{P}_{X/S}^n(\mathscr{F})
  \otimes_{\mathscr{P}_{X/S}^n}\mathscr{P}_{X/S}^n(\mathscr{G})
  \tag{16.7.8.1}
  \label{IV.16.7.8.1}
\]
(Bourbaki, \emph{Alg.}, cap.~II, 3\textsuperscript{a} ed., §~5,
n\textsuperscript{o}~1, prop.~3).

Se concluye en particular (o se ve directamente en la definición
(16.7.2.1)) que, si $\mathscr{F}$ está dotado de una estructura de
$\mathcal{O}_X$-álgebra (no necesariamente asociativa),
$\mathscr{P}_{X/S}^n(\mathscr{F})$ está canónicamente dotado de una estructura de
$\mathscr{P}_{X/S}^n$-álgebra; esta última es asociativa (resp.
conmutativa, resp. unitaria, resp. un álgebra de Lie) cuando así lo es
$\mathscr{F}$. Además, los homomorfismos canónicos
$\mathscr{P}_{X/S}^m(\mathscr{F})\to\mathscr{P}_{X/S}^n(\mathscr{F})$ para
$n\leq m$ (16.7.7) son entonces dihomomorfismos de álgebras; asimismo,
(16.7.5.1) es entonces un homomorfismo de $\mathcal{O}_X$-álgebras cuando
$\mathscr{P}_{X/S}^n(\mathscr{F})$ está dotado de su estructura de
$\mathcal{O}_X$-álgebra proveniente de su estructura de $\mathcal{O}_X$-módulo
\emph{a la derecha}.

Con las mismas notaciones, se tiene también un homomorfismo canónico de
$\mathscr{P}_{X/S}^n$-módulos
\[
  \mathscr{P}_{X/S}^n
  (\mathscr{H}\!om_{\mathcal{O}_X}(\mathscr{F},\mathscr{G}))
  \longrightarrow
  \mathscr{H}\!om_{\mathscr{P}_{X/S}^n}
  (\mathscr{P}_{X/S}^n(\mathscr{F}),\mathscr{P}_{X/S}^n(\mathscr{G}))
  \tag{16.7.8.2}
  \label{IV.16.7.8.2}
\]
(Bourbaki, \emph{Alg.}, 3\textsuperscript{a} ed., §~5,
n\textsuperscript{o}~3), que es biyectivo cuando $\mathscr{P}_{X/S}^n$ es
un $\mathcal{O}_X$-módulo \emph{localmente libre de tipo finito}
(\emph{loc. cit.}, prop.~7).
\end{env}

\begin{env}[16.7.9]
\label{IV.16.7.9}
Supongamos que se está en la situación descrita en (16.4.1); entonces, del
homomorfismo canónico $P^n(u)$ (16.4.3.3), se deduce inmediatamente un
homomorfismo canónico de $\mathcal{O}_{X'}$-bimódulos
\[
  u^*(\mathscr{P}_{X/S}^n(\mathscr{F}))
  \longrightarrow
  \mathscr{P}_{X'/S'}^n(u^*(\mathscr{F})).
  \tag{16.7.9.1}
  \label{IV.16.7.9.1}
\]
Dejamos al lector la tarea de extender a este homomorfismo las propiedades
vistas en (16.4) para el caso $\mathscr{F}=\mathcal{O}_X$.
\end{env}

\begin{remark}[16.7.10]
\label{IV.16.7.10}
La definición de $\mathscr{P}_{X/S}^n(\mathscr{F})$ en la forma (16.7.1.2)
sigue teniendo sentido cuando $\mathscr{F}$ es un haz de conjuntos cualquiera
(estando definida en (0\textsubscript{I}, 3.7.1) la imagen inversa de un haz
de conjuntos por $p_2^{(n)}$); una variante de esta definición permite
definir el «esquema de jets» (relativo a $S$) de un $X$-preesquema
cualquiera.
\end{remark}
